% answers_short.tex —— 2006 年·全国卷Ⅱ·理科 · 简短答案速查
% 由 tools/make_short_answers.py 自 answers.tex 的 \gksolhead/\gkanswer 提取，
% 并与 2006/2006_数学真题_全国2理_答案要点.md 逐题比对后生成，勿手改（改源 MD 后重跑）。

\begin{gknotice}
\textbf{使用说明}：本节只给最终答案便于快速核对，完整推导见解析版。答案为编写者独立推导，\textbf{非官方标准答案}。
\end{gknotice}

\gkpart{选择题}
\begin{center}\begin{tabular}{ccccccccccccc}
\toprule
题号 & 1 & 2 & 3 & 4 & 5 & 6 & 7 & 8 & 9 & 10 & 11 & 12 \\
\midrule
答案 & B & C & A & A & C & D & A & D & A & C & A & C \\
\bottomrule
\end{tabular}\end{center}

\gkpart{填空题}
\begin{itemize}[leftmargin=2.6em]
  \item[\textbf{13.}] $45$
  \item[\textbf{14.}] $\sqrt3$
  \item[\textbf{15.}] $\dfrac{\sqrt2}{2}$
  \item[\textbf{16.}] $25$
\end{itemize}

\gkpart{解答题}
\begin{itemize}[leftmargin=2.6em]
  \item[\textbf{17.}] (1) $\theta=-\dfrac\pi3$；(2) $|\boldsymbol a+\boldsymbol b|_{\max}=3$
  \item[\textbf{18.}] (1) $P(\xi=0)=\dfrac9{50}$，$P(\xi=1)=\dfrac{12}{25}$，$P(\xi=2)=\dfrac3{10}$，$P(\xi=3)=\dfrac1{25}$；$E\xi=\dfrac65$；(2) $\dfrac{17}{50}$
  \item[\textbf{19.}] (1) 见解析版；(2) 二面角 $A_1-AD-C_1$ 的大小为 $\dfrac\pi3$
  \item[\textbf{20.}] $a\leqslant1$。
  \item[\textbf{21.}] (1) $\overrightarrow{FM}\cdot\overrightarrow{AB}=0$；(2) $S=f(\lambda)=\dfrac12\left(\sqrt\lambda+\dfrac1{\sqrt\lambda}\right)^3$，最小值为 $4$
  \item[\textbf{22.}] (1) $a_1=\dfrac12$，$a_2=\dfrac16$；(2) $a_n=\dfrac1{n(n+1)}$
\end{itemize}
